\chapter{Methodology}

\section{Study Design}

To investigate how approachable scientific and mathematical literature is for game developers, a study was conducted.
This study, primarily a questionnaire provides quantitative statistics into the level of comprehension and the attitudes that game developers
may possess when reading, or not reading, scientific and mathematic literature, while open-ended questions provide a qualitative
perspective where uses are allowed to provide personal anecdotes of their experiences with scientific and mathematical academia.

\section{Participants}

Posts were sent out on platforms such as Discord and LinkedIn, requesting that game developers, primarily programmers fill out the questionnaire.
Due to time scheduling concerns, the survey only lasted around one week. Due to this, the participant pool was quite low; 15 unique respondents.

All the respondents met all the inclusion criteria that was established in the posts, and in the questionnaire itself.

\pagebreak

\subsection{Questionnaire}

The questionnaire (\textit{see Appendix A}) comprises:

\begin{table}[h]
    \centering
    \small
    \renewcommand{\arraystretch}{1.75}
    \caption{Questionnaire}
    \label{tab:questionnaire-short}
        \begin{tabularx}{\textwidth}{@{} l r X @{}}
        \toprule
        \textbf{Section} & \textbf{Items} & \textbf{Purpose} \\ \midrule
        Demographics & 1-3 & Determine background variables that indicate comprehension (e.g., the highest level of formal mathematics qualification) \\
        Paper Readers & 4-7 & Evaluate how participant interact, and respond to, notation, rate the clarity of lexicon on different mediums (papers vs. blogs) \\
        Barriers to Entry & 8-9 & Identify reasons for avoiding academic literature \\
        Code vs. Notation & 12-12 & Determine confidence in translating standard notation ($\sum$) into code \\ \bottomrule
    \end{tabularx}
\end{table}

\section{Data Analysis}

\subsection{Quantitative Analysis}

\subsubsection*{Demographics}

\begin{table}[H]
    \centering
    \small
    \renewcommand{\arraystretch}{1.75}
    \caption{Demographics Results}
    \label{tab:demographics-results}
        \begin{tabularx}{\textwidth}{@{} l X r r@{}}
        \toprule
        \textbf{Characteristic} & \textbf{Category} & \textbf{Count} & \textbf{Percentage} \\ \midrule
        Primary Background  & Student & 9 & 60\% \\
                            & Working in Industry & 3 & 20\% \\
                            & Indie & 2 & 13\% \\
                            & Hobbyist & 1 & 7\% \\
                            & In an adjacent field & 0 & 0\% \\ \midrule
        Maths Education Level & GCSE & 7 & 47\% \\
                              & Higher Education & 6 & 40\% \\
                              & A-Levels & 2 & 13\% \\
                              & Self-Taught & 0 & 0\% \\
                              & Other & 0 & 0\% \\ \midrule
        How often do you refer to formal or informal academic media? & Frequently & 10 & 67\% \\
        & Occasionally & 4 & 27\% \\
        & Rarely & 1 & 7\% \\
        & Never & 0 & 0\% \\
        \bottomrule
    \end{tabularx}
\end{table}

Various points can be extrapolated from the data listed in table \ref{tab:demographics-results}, such as:

\begin{itemize}
    \item The majority of respondents identified themselves as currently being students
    \item A substantial majority \textbf{(67\%)} report referring to formal or informal academic media \textit{frequently}
    \item The largest singular reported educational level is GCSE \textbf{(47\%)}. Although further education \textit{A-Levels and Higher Education} also posses a combined significant presence,
          the dominance of GCSE-educated respondents might suggest that within a large sample pool, the majority might not have completed advanced degrees.
    \item Due to the point above, it remains evident that academic content needs to be highly accessible.
\end{itemize}

\subsubsection*{Paper Readers}

\small
\renewcommand{\arraystretch}{1.2}

\begin{xltabular}{\textwidth}{@{} l X r r @{}}
    \caption{Paper Readers Results} \label{tab:paper-readers-results} \\
    \toprule
    \textbf{Question} & \textbf{Answer} & \textbf{Count} & \textbf{Percentage} \\ 
    \midrule
    \endhead
    
    \bottomrule
    \endlastfoot

    % --- Question 1 ---
    \multicolumn{4}{@{} p{\textwidth} @{}}{%
      \textbf{When reading a paper, how do you interact with mathematical notation (e.g., $\sum, \prod$)?}% % chktex 17
    } \\ \cmidrule(lr){1-4}
    & I look up their definitions on external sources & 6 & 43\% \\
    & I understand them comfortably & 4 & 29\% \\
    & I skip past the maths and try look for code snippets & 2 & 14\% \\
    & I rely on using AI to explain and give context to the notation used & 1 & 7\% \\
    & I do not understand some of the notation used & 1 & 7\% \\
    \midrule
    
    % --- Question 2 ---
    \multicolumn{4}{@{} p{\textwidth} @{}}{%
      \textbf{How would you rate the average clarity of the language used in peer-reviewed papers vs. blogs or other sources?}%
    } \\ \cmidrule(lr){1-4}
    \multicolumn{4}{@{} l @{}}{\textit{a) Peer-Reviewed Papers}} \\
    & Confusing & 8 & 57\% \\
    & Dense & 3 & 21\% \\
    & Clear & 2 & 14\% \\
    & Understandable & 1 & 7\% \\ \addlinespace
    
    \multicolumn{4}{@{} l @{}}{\textit{b) Industry blogs / other sources}} \\
    & Clear & 9 & 64\% \\
    & Understandable & 4 & 27\% \\
    & Dense & 1 & 7\% \\
    & Confusing & 0 & 0\% \\ \midrule

    % --- Question 3 ---
    \multicolumn{4}{@{} p{\textwidth} @{}}{%
      \textbf{Which of the following elements add the most value when attempting to implement a complex algorithm from a paper?}%
    } \\ \cmidrule(lr){1-4}
    & All of the above & 7 & 24\% \\
    & High level pseudo-code & 6 & 21\% \\
    & Fully implemented code snippets & 6 & 21\% \\
    & Step-by-step diagrams & 6 & 21\% \\
    & Formal writing & 4 & 14\% \\
    & None of the above & 0 & 0\% \\
    & Other & 0 & 0\% \\
    \midrule

    % --- Question 4 ---
    \multicolumn{4}{@{} p{\textwidth} @{}}{%
      \textbf{How often do you find that a paper attempts to TEACH you what a concept is, instead of just providing clarification?}%
    } \\ \cmidrule(lr){1-4}
    & Rarely & 8 & 57\% \\
    & Occasionally & 5 & 36\% \\
    & Never & 1 & 7\% \\
    & Frequently & 0 & 0\% \\

\end{xltabular}

Based the section of the survey detailed in table \ref{tab:paper-readers-results}, various key points can be extrapolated. Such as:

\begin{itemize}
    \item \textbf{Peer-reviewed papers have clarity issues}
    \begin{itemize}
        \item 78\% of respondents described peer-reviewed papers as \textbf{confusing} or \textbf{dense}; only \textbf{21\%} described them as \textbf{clear} or \textbf{understandable}
        \item In contrast, \textbf{91\%} described industry blogs and other secondary sources as \textbf{Clear}, or \textbf{Understandable}, with zero respondents finding them confusing
        Given these points, it is reasonable to suggest that readers find informal sources more approachable than formal academic papers
    \end{itemize}

    \item \textbf{Notation is a large barrier}
    \begin{itemize}
        \item Only \textbf{29\%} of respondents say that they could comfortably read mathematical notation directly within a paper
        \item \textbf{71\%} said that they required external aids, or information, such as:
        \begin{itemize}
            \item \textbf{43\%} look up definitions for notation and syntax on external sources
            \item \textbf{14\%} immediately skip all mathematics and text to find code snippets
            \item A further \textbf{14\%} said that they rely on AI tools to explain the notation, our that the outright do not understand it
        \end{itemize}

        Given these points, it is reasonable to assume that dense mathematical notation acts as a bottleneck for most readers, requiring them to leave the paper to seek context from external sources
    \end{itemize}

    \item \textbf{Readers often prefer artefacts to writing}
    
    When it comes to implementing algorithms or mathematical equations, readers prefer practical tools than formal writing:
\end{itemize}

\subsubsection*{Barriers to Entry}

\begin{xltabular}{\textwidth}{@{} l X r r @{}}
    \caption{Barriers to Entry Results} \label{tab:barriers-results} \\
    \toprule
    \textbf{Question} & \textbf{Answer} & \textbf{Count} & \textbf{Percentage} \\ 
    \midrule
    \endhead
    
    \bottomrule
    \endlastfoot

    % --- Question 1 ---
    \multicolumn{4}{@{} p{\textwidth} @{}}{%
      \textbf{What are the primary reasons that you do not use academic literature for game development?}%
    } \\ \cmidrule(lr){1-4}
    & I feel like I do not need to use academic papers & 1 & 20\% \\
    & The papers and their contents are daunting (notation, context, etc.) & 1 & 20\% \\
    & Too much academic jargon & 1 & 20\% \\
    & I feel like there is a `disconnect' from the theory in the paper and the implementation in engine & 1 & 20\% \\
    & I prefer tutorials, blogs, or other sources & 1 & 20\% \\
    & I do not know how to find relevant papers & 0 & 0\% \\
    \midrule
    
    % --- Question 2 ---
    \multicolumn{4}{@{} p{\textwidth} @{}}{%
      \textbf{If academic literature consistently used functional code samples alongside their mathematical content, how likely would you be to use them?}%
    } \\ \cmidrule(lr){1-4}
    & Neither likely nor unlikely & 1 & 100\% \\
    & Likely & 0 & 0\% \\
    & Likely & 0 & 0\% \\
\end{xltabular}

Despite only one respondent answering this section of the survey detailed in table \ref{tab:barriers-results}, various key points can be extrapolated, such as:

\begin{itemize}
    \item \textbf{Avoidance of papers is a multi-faceted problem rather than an individual issue}
    The respondent selected 5 out of the 6 listed reasons for avoiding academic literature. This might suggest that for other developers that abstain, might do so due to multiple
    reasons:
    \begin{itemize}
        \item \textbf{Lack of Relevance:} Some may perceive that academic literature simply might not be relevant for their work or use case
        \item \textbf{Comprehension:} Jargon and mathematical notation maybe considered daunting and inaccessible
        \item \textbf{Disconnect between theoretical and practical:} There may be a disconnect and difficulty translating the theoretical concepts detailed in the paper into a practical environment
        \item \textbf{Preference for alternative sources:} Alternative sources may be preferred such as tutorials, blogs, or other less formal sources
    \end{itemize}

    \item \textbf{Discovery and Access may not be a major issue} \\
    The respondent said that they do not struggle with finding or accessing papers, however this may be due to their lack of interest in doing so.
\end{itemize}

\subsubsection*{Code vs. Notation}

% --- Table 1: Individual Responses ---
\begin{table}[htbp]
  \centering
  \caption{Respondents Equation Solving Confidence}
  \label{tab:equation-solving-confidence}
  \begin{tabular}{c l l l}
    \toprule
    \textbf{Participant} & \textbf{1st Choice} & \textbf{2nd Choice} & \textbf{3rd Choice} \\
    \midrule
    1 & No idea where to start & Confident              & Somewhat confident \\
    2 & Confident              & No idea where to start & Somewhat confident \\
    3 & No idea where to start & Somewhat confident     & Confident \\
    4 & Somewhat confident     & No idea where to start & Confident \\
    5 & No idea where to start & Somewhat confident     & Confident \\
    6 & Confident              & Somewhat confident     & No idea where to start \\
    7 & Somewhat confident     & No idea where to start & Confident \\
    8 & Confident              & Somewhat confident     & No idea where to start \\
    9 & Somewhat confident     & No idea where to start & Confident \\
    \bottomrule
  \end{tabular}
\end{table}

% --- Table 2: Summary Table ---
\begin{xltabular}{\textwidth}{X r r r r r r}
  \caption{Summary of Response Choices} \label{tab:summary} \\
  \toprule
  & \multicolumn{2}{c}{\textbf{1st Choice}} & \multicolumn{2}{c}{\textbf{2nd Choice}} & \multicolumn{2}{c}{\textbf{3rd Choice}} \\
  \cmidrule(lr){2-3} \cmidrule(lr){4-5} \cmidrule(lr){6-7}
  \textbf{Confidence Level} & $count$ & \% & $count$ & \% & $count$ & \% \\
  \midrule
  \endhead

  \bottomrule
  \endlastfoot

  % --- Table Body ---
  No idea where to start & 3 & 33 & 4 & 44 & 2 & 22 \\
  Somewhat confident     & 3 & 33 & 2 & 22 & 4 & 44 \\
  Confident              & 3 & 33 & 3 & 33 & 3 & 33 \\
  \midrule
  \textbf{Total ($count$)}   & \multicolumn{2}{c}{9} & \multicolumn{2}{c}{9} & \multicolumn{2}{c}{9} \\

\end{xltabular}